\markboth{CONCLUSION AND OUTLOOK}{CONCLUSION AND OUTLOOK}
\section{Conclusion and Outlook}\label{sec:conclusion}

\subsection{Summary}\label{sec:conclusion:summary}

This paper presented Training Copilot, an open-source, multi-agent sports analytics platform that unifies fragmented fitness data behind a single conversational interface, addressing data fragmentation across wearable ecosystems, the lack of cross-domain contextual reasoning, and the rigid architectures of current platforms. The MCP-based integration layer achieves single-line extensibility, verified empirically by adding the Telegram proxy and flythrough servers during development. The LangGraph-based multi-agent system coordinates six domain-scoped specialists via A2A, with parallel orchestration keeping latency manageable and scoped ToolHosts reducing tool misuse; the recording-and-clipping pattern decouples data fidelity from LLM context constraints, a RAG-based fitness specialist grounds advice in verifiable published sources with explicit citation preservation, and a coach specialist turns a structured race goal into a guardrail-validated multi-week training plan whose arithmetic is deterministic code rather than model output. The automated evaluation framework combines LLM judges with a deterministic trace-based grounding scorer for reproducible quality assessment.

\subsection{Future Work}\label{sec:conclusion:next}

\textbf{Field study.} The most important next step is a controlled study with real athletes: recruiting 10-20 participants from the KIT sports community, deploying Training Copilot as their training assistant for four weeks, and measuring both quantitative metrics (training consistency, load progression) and qualitative feedback (perceived usefulness, trust). The automated evaluation (Section~\ref{sec:evaluation}) assesses conversational quality; only longitudinal human evaluation can measure impact on actual training decisions.

\textbf{Knowledge base expansion.} The RAG corpus comprises five contemporary German-language sports-science textbooks; incorporating open-access publications (e.g.\ from PubMed Central) would broaden its coverage and let the fitness specialist cite current research in more languages. Training Copilot now generates multi-week periodised plans~\citep{grgic_periodization_2017} through the coach specialist and the athlete server's deterministic scaffold (Section~\ref{sec:data:athlete}); the scaffold prescribes phases, weekly volumes, and the long-run line, while the LLM chooses and explains individual workouts. Extending this toward richer agent-authored plan structures that stay continuously adjusted against live load and calendar is a natural next step that reuses the existing agent layer.

\textbf{Multi-tenant deployment.} Production deployment requires per-user credential vaults (replacing the current plaintext \texttt{.tokens/} directory), session-scoped ToolHost instances, and the security measures outlined in Section~\ref{sec:discussion:limitations}. The stateless architecture (no server holds per-request state) supports horizontal scaling; the orchestrator's in-process trace assembly is the main bottleneck for this. A field study (above) needs this groundwork too: handing out accounts and devices to participants requires proper account management, which the current single-user \texttt{.tokens/} setup does not provide.

\textbf{Ecosystem evolution.} Strava's announcement of an official MCP server~\citep{strava_api_update_2026} supports the architecture's bet on protocol-based integration and suggests first-party MCP servers from data providers will gradually replace third-party API wrappers. As MCP adoption grows, community-contributed servers (nutrition, air quality, social rides) could be added with a single URL; enabling token-level SSE streaming and caching frequently repeated tool calls would further reduce the 8-15 second latency discussed in Section~\ref{sec:discussion:limitations}.

\textbf{Hallucination mitigation and answer verification.} Training Copilot reduces hallucination mainly \textit{by construction}: it grounds every answer in tool results or retrieved passages, iterates over real observations in the ReAct loop rather than free-generating~\citep{yao_react_2023}, and preserves citations. What it does not yet do is \textit{verify} that a synthesised answer stays faithful to that evidence. Because the trace already records every tool call and retrieved passage (Section~\ref{sec:agents:recording}), a post-synthesis verification stage can be added without changing the specialists: a lightweight critic step (or self-consistency across sampled answers) that flags any claim unsupported by the trace, the retrieval-faithfulness scorer planned as the sixth evaluation dimension (Section~\ref{sec:evaluation:scorers}), and explicit abstention when the evidence is thin. We regard systematic answer verification as the most valuable next step toward trustworthy, reliable coaching advice, and are actively extending the evaluation harness in this direction.

\textbf{Experience-aware route planning.} The route specialist already enriches a route with points of interest along the way: given a planned loop or point-to-point route, it finds a caf\'e, bakery, or water fountain that lies on the track (Section~\ref{sec:data:other}). Two things remain. First, the circular-route generator is a proof of concept: it produces a loop of a target distance from a start point, but other route shapes (a figure-eight, an out-and-back along a river, a landmark tour) would need further work, which we defer to stay focused on the core analytics and planning path. Second, a request like ``plan a scenic bike loop where I can swim in a lake around halfway and stop for ice cream on the way back'' needs the specialist to select and sequence several \textit{different} stops at specific points of the route, not just one; combined with a higher-quality curated route source (Section~\ref{sec:discussion:limitations}), this would match consumer route planners on experience while keeping the multi-source agentic reasoning that sets Training Copilot apart.

\subsection{Closing Remarks}\label{sec:conclusion:closing}

Training Copilot shows that protocol-based data integration (MCP), multi-agent coordination (A2A), and domain-scoped specialist agents (LangGraph ReAct) can bridge the gap between fragmented fitness data and practical training advice. The system runs in real time on commodity hardware and can be extended by anyone who can write a Python function and a one-line configuration entry. As standardised agent protocols mature and wearable ecosystems become more interoperable, architectures like Training Copilot's, where the intelligence sits in the coordination layer rather than the data layer, offer a viable pattern for agentic applications beyond sports analytics.
